\emph{প্রমাণপদ্ধতি} দিয়ে বাক্যগুলির মধ্যে অনুসিদ্ধান্ত ও
সামঞ্জস্যকে বিশুদ্ধভাবে সংকেতবিন্যাসের ভিত্তিতে চরিত্রায়িত করা
যায়। \emph{স্বাভাবিক নিষ্পাদন} এমনই একটি প্রমাণপদ্ধতি।
এতে একটি \emph{নিষ্পাদন} সূত্রগুলির একটি বৃক্ষ নিয়ে গঠিত।
!!a{derivation}-এর একেবারে উপরের সূত্রগুলি হলো \emph{অনুমিতি}।
!!{derivation} সঠিক হতে হলে বাকি প্রত্যেকটি সূত্রকে কয়েকটি
\emph{অনুমানবিধির} একটির সাহায্যে যথাযথভাবে প্রতিষ্ঠা করতে হয়।
এই বিধিগুলি জোড়ায় জোড়ায় আসে। প্রতিটি সংযোজক ও পরিমাণসূচকের
জন্য একটি প্রবর্তন-বিধি ও একটি অপসারণ-বিধি থাকে। যেমন, সূত্র~$!A$
একটি $\Elim{\lif}$ বিধি দিয়ে প্রতিষ্ঠা করা হলে আগের সূত্রগুলি
(\emph{পূর্বধারণাগুলি}) অবশ্যই $!B \lif !A$ ও $!B$ হবে
(কোনো একটি~$!B$-এর জন্য)। কিছু অনুমানবিধিতে অনুমিতিগুলিকে
\emph{অবমুক্ত} করাও যায়। যেমন, $!A \lif !B$ যদি $!B$ থেকে
$\Intro{\lif}$ প্রয়োগ করে অনুমান করা হয়, তবে পূর্বধারণা~$!B$-তে পৌঁছানোর
!!{derivation}-এ~$!A$ যত জায়গায় অনুমিতি হিসেবে এসেছে,
সেগুলি অবমুক্ত করা যেতে পারে। সেগুলিকে একটি লেবেল দেওয়া হয়,
যা অনুমানের ধাপটিতেও লেখা হয়।

অন্তিম সূত্র~$!A$-সহ একটি !!a{derivation} থাকলে এবং তার সব
অনুমিতি অবমুক্ত হলে বলি $!A$ একটি \emph{উপপাদ্য} এবং লিখি~$\Proves !A$।
অন্তিম সূত্র~$!A$-সহ কোনো নিষ্পাদনের সব অনবমুক্ত অনুমিতি সেট~$\Gamma$-তে
থাকলে বলি $!A$ হলো
~$\Gamma$ \emph{থেকে নিষ্পাদনযোগ্য} এবং লিখি $\Gamma \Proves !A$।
$\Gamma \Proves \lfalse$ হলে বলি $\Gamma$ \emph{অসঙ্গত}, অন্যথায়
\emph{সঙ্গত}। এই ধারণাগুলি পরস্পর সম্পর্কিত। যেমন,
$\Gamma \Proves !A$ হয় ঠিক তখনই, যখন $\Gamma \cup \{\lnot !A\}$
অসঙ্গত। অনুরূপ অর্থগত ধারণাগুলির সঙ্গেও এদের সম্পর্ক আছে। যেমন,
$\Gamma \Proves !A$ হলে $\Gamma \Entails !A$। প্রমাণপদ্ধতির এই
ধর্মকে \emph{বিশুদ্ধতা} বলে—$\Gamma$ থেকে যা নিষ্পাদন করা যায়,
তা~$\Gamma$ থেকে অবশ্যই অনুসৃত হয়। !!{derivation}s-এর দৈর্ঘ্যের উপর
আরোহ করে \emph{বিশুদ্ধতা উপপাদ্য} প্রমাণ করা হয়। তাতে দেখানো হয়
যে প্রতিটি অনুমানে পূর্বধারণাগুলি তাদের অনবমুক্ত অনুমিতিগুলি থেকে
অনুসৃত হলে সিদ্ধান্তটিও তার খোলা অনুমিতিগুলি থেকে অনুসৃত হয়।
